\usepackage{hyperref}
\usepackage{xcolor}
\usepackage{calc}
\usepackage{graphicx}
\usepackage{tikz}
\usepackage{fontspec}
\usepackage{fontawesome5}
\usepackage{titlesec}
\usepackage{enumitem}
\usepackage{fancybox}
\usepackage{xeCJK}

\hypersetup{hidelinks}
\urlstyle{same}

%%%%%%%%%%%%%%%%%%%%
% 设置
%%%%%%%%%%%%%%%%%%%%

\providecommand{\ResumeItemSep}{0.12em}
\providecommand{\ResumeTopSep}{0.18em}
\providecommand{\ResumeArrayStretch}{1.12}
\providecommand{\ResumeLineSpread}{1.08}
\providecommand{\ResumeSectionBefore}{0.55em}
\providecommand{\ResumeSectionAfter}{0.28em}

\setlength{\parindent}{0pt}
\setlength{\parskip}{0pt}
\pagenumbering{gobble}
\setlist[itemize]{
    leftmargin=1.35em,
    itemsep=\ResumeItemSep,
    topsep=\ResumeTopSep,
    parsep=0pt,
    partopsep=0pt
}
\setlist[enumerate]{leftmargin=*}
\renewcommand{\arraystretch}{\ResumeArrayStretch}
\linespread{\ResumeLineSpread}

\titleformat{\section}
  {\Large\bfseries\raggedright}
  {}{0em}
  {}
  [{\color{secondarycolor}\titlerule}]
\titlespacing*{\section}{0pt}{\ResumeSectionBefore}{\ResumeSectionAfter}

\usepackage[
    a4paper,
    left=1.1cm,
    right=1.1cm,
    top=0.75cm,
    bottom=0.85cm,
    nohead
]{geometry}

% 字体设置
% 说明：Noto Serif SC 没有原生斜体字型。将 CJK 的 italic 字型映射到常规字形
%（与过去“找不到斜体、默认替换为常规字形”的实际渲染一致），可消除
% “Font shape `TU/NotoSerifSC(0)/m/it' undefined” 告警；如需真正斜体观感，
% 可在此行追加 ItalicFeatures={FakeSlant=0.18}。
\setCJKmainfont[
    Path=fonts/,
    Extension=.otf,
    BoldFont=*-Bold,
    ItalicFont={NotoSerifSC},
]{NotoSerifSC}

% 自定义颜色
\definecolor{primarycolor}{RGB}{200, 60, 60}
\definecolor{secondarycolor}{RGB}{200, 68, 28}

\newlength{\iconwidth}
\setlength{\iconwidth}{1.5em}

\newcommand{\sectionicon}[2]{\makebox[\iconwidth][c]{\color{primarycolor}{#1}}\quad #2}

\newcommand{\projectheading}[3]{%
    \noindent\makebox[\textwidth][s]{%
        \makebox[0pt][l]{{\large \textbf{#1}}}%
        \hfill
        \makebox[0pt][c]{{\normalsize \textbf{#2}}}%
        \hfill
        \makebox[0pt][r]{#3}%
    }%
}
